\scanpage{53}
\margin{\handwritten{\unsure{suite de|mieux de} pentes de Hodge}}
% slanted: leader runs from the marginal words to the first primed lambda.
\[
\begin{aligned}
&\handwritten{\left\{\begin{aligned}
\sum_{1}^{i}\lambda'_j&\leq\deleted{\unsure{\vee}}\sum_{1}^{i}\lambda_j\\
\sum_{\unsure{n|N|1}}^{\deleted{\unsure{n}}}\lambda'_j&=\sum\lambda_j
\end{aligned}\right.}\\[-1pt]
&\qquad\handwritten{\overset{\unsure{''|=}}{\sum}\,i h_i^{\prime}}
\end{aligned}
\]

\noindent\rule{\linewidth}{0.3pt}\par

\deleted{\handwritten{\uline{\unsure{Sur un th. de Tate.}}}} Sur un théorème de Tate. (Devrait venir avec foncteur mystérieux !)

C'est le théorème suivant (Driebergen ....):

\uline{Th.} (Tate): Soit $S$ un schéma noethérien normal intègre dont le corps des
fonctions \inserted{$k(\handwritten{\eta})$} est de car. zero, Alors le foncteur
\[
\mathrm{BT}(S)\handwritten{\longrightarrow}\mathrm{BT}(\handwritten{\eta})
\]
est pleinement fidèle.

Ainsi, un groupe de BT sur $S$ est c\deleted{\unsure{b}}onnu \inserted{à isom unique près} quand on connait la représentation $p$-adique qu'\deleted{elle}il définit de $\handwritten{\pi}_{1}(\handwritten{\eta})=\mathrm{Gal}(\overline{k(\handwritten{\eta})}/k(\handwritten{\eta}))$. Prenant d'abord
le cas où $S$ est le spectre d'un anneau de valuation discrète complet $V$
d'inégales car. , de corps \deleted{des} résiduel $k$ parfait, la question se pose donc

a) De caractériser les modules galoisiens qu'on peut obtenir (Tate donne
des conditions nécessaires: les modules en question doivent être du type de
"Hodge-Tate" );

b) De trouver un foncteur "explicite" sur la catégorie de ces modules galoisiens, vers les F-cristaux sur $k\quad M$ muni d'une filtration de $M\otimes_W\deleted{\unsure{V}}L$.

\handwritten{\uline{Problème}} Le théorème de Tate suggère de plus que le "foncteur mystérieux"
qui va des F-cristaux sur $k$ \deleted{\unsure{m}}filtrés sur $L$, vers les $\underline Q_p$-modules galoisiens
sur $L$, est pleine\deleted{\unsure{n}}ment fidèle, et qu'on doit pouvoir donc \deleted{définir un foncteu}
expliciter un foncteur quasi-\deleted{\unsure{c}}inverse (non moins mystérieux), allant d'une
sous-catégorie pleine convenable des $\underline Q_p$-modules \deleted{\unsure{G}}galoisiens sur $L$ (sans
doute contenue dans celle des modules galoisiens du type de Hodge-Tate)
vers les F-cristaux sur $k$ filtrés sur $L$.

D'autre part, on conjec\deleted{\unsure{y}}ture que dans le th. de Tate ci-dessus, la
restriction au cas où \deleted{$S$}\inserted{$k(\handwritten{\eta})\handwritten{=L}$} est de car. nulle est inutile, Le seul cas qui offre
% archive sheet number: 52

